\subsection{Radial mass and thin bands}\label{sec:radial-mass}
In the next two lemmas we turn bounds on $J_N$ at a few radii into bounds
on root counts, and in Section~\ref{sec:concentration} these bounds are
supplied by Theorem~\ref{thm:sign-logarithms} outside exceptional events
with summable probabilities.

\begin{lemma}[Annular mass]\label{lem:annular-mass}
Let $K>0$, $0<\eta<1$ and $N\ge2K$, let $f_N$ be a nonzero polynomial of
degree at most $N$, and let $r_i=1+x_i/N$ with
$(x_1,x_2,x_3,x_4)=(-K,-(1-\eta)K,(1-\eta)K,K)$.
Let $\varepsilon,D\ge0$ be such that
\begin{align*}
 \left|J_N(r_i)-\log\sigma_N(r_i)+\gamma/2\right|&\le\varepsilon,\\
 \left|\log\sigma_N(r_i)-\tfrac12\log N-F(x_i)\right|&\le D/N,
 \qquad i=1,2,3,4.
\end{align*}
Then
\begin{align*}
 &\frac{\#\{\alpha:f_N(\alpha)=0,\ ||\alpha|-1|>K/N\}}N\\
 &\qquad\le\frac2{\eta K}\{F(-(1-\eta)K)-F(-K)\}
   +\frac4{\eta K}\Big(\varepsilon+\frac DN\Big)\\
 &\qquad\le\frac1{\eta K}\log\frac1{1-\eta}
   +\frac4{\eta K}\Big(\varepsilon+\frac DN\Big).
\end{align*}
In particular, for $\eta=1/2$,
\begin{align*}
 \frac{\#\{\alpha:f_N(\alpha)=0,\ ||\alpha|-1|>K/N\}}N
 &\le\frac2K\log\frac2{1+e^{-K}}+\frac8K\Big(\varepsilon+\frac DN\Big)\\
 &\le\frac{2\log2}K+\frac8K\Big(\varepsilon+\frac DN\Big).
\end{align*}
\end{lemma}

By Lemma~\ref{lem:profile-rate}, the second condition holds with
$D=K^2+e^{4K}/2$.

\begin{proof}
We divide the two secants of Lemma~\ref{thm:T2} by the linear spacing
$\eta K$ in place of the logarithmic one, replace $J_N$ by
$\frac12\log N+F$, and evaluate the leading term by
Proposition~\ref{prop:profiles}.
Let $m$ be the number of roots with $||\alpha|-1|>K/N$.
Since $0<1+x/N\le1$ for $-K\le x\le0$, because $N>K$, and
$1+x/N\ge1$ for $x\ge0$, the denominators satisfy
\begin{align*}
 N\log(r_2/r_1)&=\int_{-K}^{-(1-\eta)K}\frac{\mathrm dx}{1+x/N}\ge\eta K,\\
 N\log(r_4/r_3)&=\int_{(1-\eta)K}^{K}\frac{\mathrm dx}{1+x/N}\le\eta K.
\end{align*}
By \eqref{eq:jensen-integral}, $J_N(r)$ is nondecreasing in $r$, so that
the numerators $J_N(r_2)-J_N(r_1)$ and $J_N(r_4)-J_N(r_3)$ are
nonnegative. By the two hypotheses, in which $\gamma/2$ and
$\frac12\log N$ cancel in differences,
\begin{align*}
 J_N(r_2)-J_N(r_1)&\le F(x_2)-F(x_1)+2\Big(\varepsilon+\frac DN\Big),\\
 J_N(r_4)-J_N(r_3)&\ge F(x_4)-F(x_3)-2\Big(\varepsilon+\frac DN\Big).
\end{align*}
By the second inequality of the bound on $m/N$ in the proof of
Lemma~\ref{thm:T2}, $\deg f_N\le N$, the bounds on the denominators and
the two bounds above, we have
\begin{align*}
 \frac mN
 &\le1+\frac{J_N(r_2)-J_N(r_1)}{\eta K}-\frac{J_N(r_4)-J_N(r_3)}{\eta K}\\
 &\le1+\frac1{\eta K}\{F(x_2)-F(x_1)-F(x_4)+F(x_3)\}
 +\frac4{\eta K}\Big(\varepsilon+\frac DN\Big)\\
 &=\frac2{\eta K}\{F(-(1-\eta)K)-F(-K)\}
 +\frac4{\eta K}\Big(\varepsilon+\frac DN\Big),
\end{align*}
where the last line is the evaluation of \eqref{eq:T2-limit} in the proof
of Proposition~\ref{prop:profiles}, and \eqref{eq:profile-eta-bound}
bounds its first term.
\end{proof}

The two secants contribute equally to the leading term
$\frac2{\eta K}\{F(-(1-\eta)K)-F(-K)\}$ of
Lemma~\ref{lem:annular-mass}: the inner secant of Lemma~\ref{thm:T2}
controls the roots in $|z|<1-K/N$, and the outer secant controls those in
$|z|>1+K/N$.
By Lemma~\ref{lem:profile-rate} and $N\log(r_2/r_1)\to\eta K$, the inner
secant $(\log\sigma_N(r_2)-\log\sigma_N(r_1))/(N\log(r_2/r_1))$ tends to
\[
 \frac{F(-(1-\eta)K)-F(-K)}{\eta K}.
\]
There are at most $N-\nu_N(r_4)$ roots in $|z|>1+K/N$, and
$\nu_N(r_4)/N$ is at least the outer secant, as in the proof of
Lemma~\ref{thm:T2}. By the same limits and the symmetry $F(x)=x+F(-x)$,
$1$ minus the outer secant tends to
\begin{align*}
 1-\frac{F(K)-F((1-\eta)K)}{\eta K}
 &=1-\frac{\eta K+F(-K)-F(-(1-\eta)K)}{\eta K}\\
 &=\frac{F(-(1-\eta)K)-F(-K)}{\eta K}.
\end{align*}
As $\eta\to0$ the leading term tends to the fraction of roots outside the
annulus given by Theorem~\ref{thm:radial-profile}, since $F'=\Phi$ and
$\Phi(x)+\Phi(-x)=1$:
\[
 \lim_{\eta\to0}\frac2{\eta K}\{F(-(1-\eta)K)-F(-K)\}
 =2\Phi(-K)=1-\{\Phi(K)-\Phi(-K)\}.
\]
Its upper bound $\log(1/(1-\eta))/(\eta K)$ tends to $1/K$, and the value
$\eta=1/2$ gives $2\log2/K$, which suffices when only $K\to\infty$
matters.

\begin{lemma}[Thin bands]\label{lem:thin-band}
Fix $0<\kappa\le1/64$, let $N\ge4$, and let $f_N$
be a nonzero polynomial of degree at most $N$.
Suppose that for some $\varepsilon\ge0$,
\[
 \left|J_N(1+qN^{-1-\kappa})-\log\sigma_N(1+qN^{-1-\kappa})+\gamma/2\right|\le\varepsilon,
 \qquad q=-2,-1,0,1,2.
\]
Then, for $q=-1,0,1$,
\[
 \left|\frac{\nu_N(1+qN^{-1-\kappa})}N-\frac12\right|
 \le2N^{-\kappa}+4\varepsilon N^{\kappa}.
\]
\end{lemma}

\begin{proof}
We follow the proof of Lemma~\ref{thm:T2}, now at five radii spaced by
$N^{-1-\kappa}$. At each of the three middle radii, the count $\nu_N$ lies
between the slopes of two secants of $t\mapsto J_N(e^t)$, and a bound on
the second derivative of $\log\sigma_N(e^t)$ shows that the slopes of the
corresponding secants of $\log\sigma_N(e^t)$, divided by $N$, are within
$2N^{-\kappa}$ of $1/2$.
Let $g(t)=\log\sigma_N(e^t)$.
Differentiating the finite sum, we see that
\begin{align*}
 g'(t)&=\frac{\sum_{k=0}^Nke^{2kt}}{\sum_{k=0}^Ne^{2kt}},\\
 g''(t)&=2\left\{
 \frac{\sum_{k=0}^Nk^2e^{2kt}}{\sum_{k=0}^Ne^{2kt}}
 -\left(\frac{\sum_{k=0}^Nke^{2kt}}{\sum_{k=0}^Ne^{2kt}}\right)^2
 \right\}.
\end{align*}
That is, under the probability weights proportional to $e^{2kt}$ on
$\{0,1,\dots,N\}$, $g'(t)$ is the mean of $k$ and $g''(t)$ is twice its
variance $\operatorname{Var}_t(k)$. At $t=0$ the weights are uniform, so
that $g'(0)=\frac1{N+1}\sum_{k=0}^Nk=\frac N2$.
For a random variable $Y$ with values in $[m,M]$, since
$\mathbb E(Y-m)(M-Y)\ge0$, we have
\begin{align*}
 \operatorname{Var}(Y)
 &=(\mathbb EY-m)(M-\mathbb EY)-\mathbb E(Y-m)(M-Y)\\
 &\le(\mathbb EY-m)(M-\mathbb EY)
 \le\frac{(M-m)^2}{4}.
\end{align*}
Applying this bound with $Y=k$, $m=0$ and $M=N$, we obtain
\begin{equation}
 g''(t)=2\operatorname{Var}_t(k)\le2\cdot\frac{N^2}4=\frac{N^2}2.
                                                        \label{eq:layer2-gpp}
\end{equation}
Let $t_q=\log(1+qN^{-1-\kappa})$ for $q=-2,-1,0,1,2$.
Since $N\ge4$, we have $N^{-1-\kappa}\le1/4$, and since
$|\log(1+y)|\le|y|/(1-|y|)\le2|y|$ for $|y|\le1/2$, the logarithmic radii
and their gaps satisfy
\begin{align*}
 |t_q|&\le2|q|N^{-1-\kappa}\le4N^{-1-\kappa},\\
 t_{q+1}-t_q
 &=\int_{1+qN^{-1-\kappa}}^{1+(q+1)N^{-1-\kappa}}\frac{\mathrm du}{u}
 \ge\frac{N^{-1-\kappa}}{1+(q+1)N^{-1-\kappa}}
 \ge\frac{N^{-1-\kappa}}2,\qquad q=-2,-1,0,1.
\end{align*}
By the hypothesis at $t_{-2},\dots,t_2$, in which $\gamma/2$ cancels in
differences, by \eqref{eq:jensen-integral} and since $\nu_N(e^t)$ is
nondecreasing in $t$, for $q=-1,0,1$ we have
\begin{align*}
 \frac{g(t_q)-g(t_{q-1})-2\varepsilon}{N(t_q-t_{q-1})}
 &\le\frac{J_N(e^{t_q})-J_N(e^{t_{q-1}})}{N(t_q-t_{q-1})}
 =\frac1{N(t_q-t_{q-1})}\int_{t_{q-1}}^{t_q}\nu_N(e^t)\dd t\\
 &\le\frac{\nu_N(1+qN^{-1-\kappa})}N
 \le\frac1{N(t_{q+1}-t_q)}\int_{t_q}^{t_{q+1}}\nu_N(e^t)\dd t\\
 &=\frac{J_N(e^{t_{q+1}})-J_N(e^{t_q})}{N(t_{q+1}-t_q)}
 \le\frac{g(t_{q+1})-g(t_q)+2\varepsilon}{N(t_{q+1}-t_q)}.
\end{align*}
Let $[a,b]$ be $[t_{q-1},t_q]$ or $[t_q,t_{q+1}]$. Since $g'(0)=N/2$, and
$|g'(t)-g'(0)|\le\frac{N^2}2|t|$ by \eqref{eq:layer2-gpp} and the mean
value theorem, with $|t|\le4N^{-1-\kappa}$ on $[a,b]$, we have
\begin{align*}
 \left|\frac{g(b)-g(a)}{N(b-a)}-\frac12\right|
 &=\frac1{N(b-a)}\left|\int_a^b(g'(t)-g'(0))\dd t\right|\\
 &\le\frac1{N(b-a)}\int_a^b|g'(t)-g'(0)|\dd t
 \le\frac1{N(b-a)}\int_a^b\frac{N^2}2|t|\dd t\\
 &\le\frac N2\cdot4N^{-1-\kappa}
 =2N^{-\kappa}.
\end{align*}
Moreover, since $b-a\ge N^{-1-\kappa}/2$, we have
$2\varepsilon/(N(b-a))\le4\varepsilon N^\kappa$. By the secant bounds and
the last two estimates, we have
\begin{align*}
 \frac{\nu_N(1+qN^{-1-\kappa})}N-\frac12
 &\le\left(\frac{g(t_{q+1})-g(t_q)}{N(t_{q+1}-t_q)}-\frac12\right)
   +\frac{2\varepsilon}{N(t_{q+1}-t_q)}
 \le2N^{-\kappa}+4\varepsilon N^{\kappa},\\
 \frac12-\frac{\nu_N(1+qN^{-1-\kappa})}N
 &\le\left(\frac12-\frac{g(t_q)-g(t_{q-1})}{N(t_q-t_{q-1})}\right)
   +\frac{2\varepsilon}{N(t_q-t_{q-1})}
 \le2N^{-\kappa}+4\varepsilon N^{\kappa}.
\end{align*}
\end{proof}
